% titel: Blatt 2, A4: Stoffmenge Ethanol aus Volumen und Dichte
% Teil II: der durchgerechnete Weg.

\subsection*{Gegeben und gesucht}
Volumen Ethanol $V = \qty{30}{mL}$, Dichte $\varrho =
\qty{0.789}{g/mL}$, Summenformel $\mathrm{C_2H_6O}$. Gesucht die Stoffmenge
$n$ an Ethanol und die Stoffmengen an einzelnen Atomen ($\mathrm{H,C,O}$).

\subsection*{Masse}
\begin{equation}
  m = \varrho \cdot V = \qty{0.789}{g/mL} \cdot \qty{30}{mL}
    = \qty{23.67}{g}.
\end{equation}

\subsection*{Stoffmenge Ethanol}
Molmasse $\mathrm{C_2H_6O}$:
$M = 2\cdot 12.011 + 6\cdot 1.008 + 15.999 = 46.07\ \mathrm{g/mol}$.
\begin{equation}
  n(\mathrm{C_2H_6O}) = \frac{m}{M}
    = \frac{23.67}{46.07} = 0.514\ \mathrm{mol}.
\end{equation}

\subsection*{Stoffmenge der einzelnen Atome}
Pro Molekül: $2$ C, $6$ H, $1$ O.
\begin{align*}
  n(\mathrm{C}) &= 2\cdot 0.514 = 1.028\ \mathrm{mol},\\
  n(\mathrm{H}) &= 6\cdot 0.514 = 3.084\ \mathrm{mol},\\
  n(\mathrm{O}) &= 1\cdot 0.514 = 0.514\ \mathrm{mol}.
\end{align*}

\begin{ergebnis}
\begin{align*}
  n(\mathrm{C_2H_6O}) &= \qty{0.514}{mol},\\
  n(\mathrm{C}) &= \qty{1.03}{mol},\quad
  n(\mathrm{H}) = \qty{3.08}{mol},\quad
  n(\mathrm{O}) = \qty{0.51}{mol}.
\end{align*}
\end{ergebnis}

\begin{kontrolle}
Die Atomstoffmengen stehen im Verhaeltnis $1.03:3.08:0.51 = 2:6:1$, exakt die
Indizes der Summenformel $\mathrm{C_2H_6O}$. Die Summe
$n(\mathrm{C})+n(\mathrm{H})+n(\mathrm{O}) = 4.63\ \mathrm{mol}$ ist größer
als $n(\mathrm{C_2H_6O})$, weil jedes Molekül mehrere Atome traegt.
\end{kontrolle}

\begin{merksatz}
Volumen über die Dichte in Masse umrechnen, dann $n = m/M$. Für die
Atomsorten die Stoffmenge des Moleküls mit den Summenformel-Indizes
multiplizieren.
\end{merksatz}
